% =====================================================================
% Relatorio da Entrega 2 -- Computacao Distribuida (PUC Minas)
% Mesmo estilo do relatorio da Entrega 1 (artigo SBC/JBCS, duas colunas).
% Compilar: pdflatex relatorio.tex (2x).
% =====================================================================
\documentclass[10pt,twocolumn]{article}

\usepackage[utf8]{inputenc}
\usepackage[T1]{fontenc}
\usepackage{times}
\usepackage[a4paper,top=2.3cm,bottom=2.3cm,left=2.0cm,right=2.0cm,columnsep=0.8cm]{geometry}
\usepackage{graphicx}
\usepackage{url}
\usepackage[brazil]{babel}
\usepackage[font=small,labelfont=bf]{caption}

\setlength{\parindent}{1em}
\setlength{\parskip}{1pt}

% Marcador para valores a medir/confirmar pelo grupo antes da entrega
\newcommand{\medir}[1]{\textbf{[#1]}}

\newcommand{\autor}[2]{\noindent\textbf{#1}\ {\small[ Universidade PUC MG $|$ \texttt{#2} ]}\par}

\sloppy

\begin{document}

\twocolumn[{%
\begin{flushleft}
{\LARGE\bfseries Banco de Dados Distribuído com Replicação Síncrona,\\
Two-Phase Commit e Eleição de Líder\par}
\vspace{6pt}
\autor{Edson Pimenta de Almeida}{1435541@sga.pucminas.br}
\autor{Felipe Carvalho de Paula Silva}{1319360@sga.pucminas.br}
\autor{João Lucas de Melo Quintão}{1165769@sga.pucminas.br}
\autor{Luís Augusto Starling Toledo}{1323060@sga.pucminas.br}
\autor{Juan Pablo Ramos de Oliveira}{1467614@sga.pucminas.br}
\autor{Luiz Gabriel Milione Assis}{1493412@sga.pucminas.br}
\autor{Túlio Gomes Braga}{1441272@sga.pucminas.br}
\vspace{6pt}

{\small\textbf{Resumo.} Este trabalho apresenta um banco de dados
chave--valor distribuído desenvolvido em Python com gRPC. As escritas são
coordenadas por um líder eleito dinamicamente pelo algoritmo do valentão e
aplicadas de forma atômica em todas as réplicas vivas por meio do protocolo
\emph{Two-Phase Commit} (2PC), com exclusão mútua distribuída por bloqueio
de chaves, quórum majoritário e relógios lógicos de Lamport. O sistema
tolera falhas em várias camadas: detecção por \emph{heartbeat}, reeleição
automática de líder, registro de escrita antecipada (WAL) para recuperação
após quedas, transferência de estado para nós que reingressam e um
protocolo de terminação que resolve transações pendentes quando o
coordenador falha. A avaliação combina testes automatizados, injeção de
falhas e um teste de estresse com clientes concorrentes, verificando ao
final a consistência entre os nós.\par}
\vspace{3pt}
{\small\textbf{Palavras-chave:} Banco de Dados Distribuído, gRPC,
Two-Phase Commit, Eleição de Líder, Tolerância a Falhas, Exclusão Mútua,
Relógio Lógico\par}
\vspace{8pt}
\end{flushleft}
}]

\section{Introdução}

Um banco de dados distribuído precisa manter os mesmos dados em várias
máquinas ao mesmo tempo. Esse objetivo esconde três problemas clássicos de
sistemas distribuídos: (i)~uma escrita deve se tornar visível em
\emph{todas} as réplicas ou em \emph{nenhuma}, mesmo que parte delas falhe
no meio da operação; (ii)~leituras atendidas por réplicas diferentes devem
devolver o mesmo valor; e (iii)~a queda de qualquer máquina, inclusive a
que coordena as escritas, não pode interromper o serviço nem corromper os
dados.

O sistema desenvolvido é um banco chave--valor replicado, acessado por um
cliente de linha de comando com operações de escrita (\texttt{put},
\texttt{update}, \texttt{delete}) e de leitura (\texttt{get},
\texttt{exists}, \texttt{list}, \texttt{size}, \texttt{getAll}). As
principais decisões metodológicas foram: replicação \emph{síncrona}
coordenada por um líder, priorizando consistência sobre disponibilidade;
transações distribuídas com 2PC complementado por um protocolo de
terminação; eleição de líder pelo algoritmo do valentão com salvaguarda
contra estado defasado; e separação estrita entre a lógica distribuída e o
transporte de rede, para permitir testes determinísticos de cada algoritmo.

\section{Tecnologias Utilizadas}

O sistema foi desenvolvido em Python~3, com gRPC para comunicação remota e
Protocol Buffers para definição do contrato de serviço. O ambiente é
isolado com \texttt{venv} e as dependências são fixadas em
\texttt{requirements.txt}. Os algoritmos distribuídos (relógio lógico,
exclusão mútua, eleição, 2PC, detector de falhas, terminação e quórum)
foram implementados manualmente pelo grupo; a biblioteca gRPC é usada
apenas como transporte.

\section{Arquitetura do Sistema}

O sistema é formado por $N$ nós idênticos (os experimentos usam $N=3$, mas
a topologia é configurável por linha de comando) e por clientes. Todos os
nós executam o mesmo programa e os papéis são definidos em tempo de
execução:

\begin{itemize}\itemsep2pt
  \item \textbf{Líder}: o nó vivo de maior identificador; é o único que
    coordena escritas, executando o 2PC sobre as réplicas vivas, e é a
    fonte de estado para nós que reingressam.
  \item \textbf{Réplicas (participantes)}: armazenam os dados, votam no
    2PC, atendem leituras e vigiam o líder por \emph{heartbeat}; se ele
    morre, disparam uma eleição.
  \item \textbf{Cliente}: conhece a lista de nós; envia escritas ao líder
    (se contatar o nó errado recebe \texttt{NOT\_LEADER} com o endereço
    correto e se redireciona) e lê de qualquer nó.
\end{itemize}

A topologia lógica é uma \textbf{malha completa}: todo nó se comunica
diretamente com todos os outros. O \emph{heartbeat} é todos-para-todos a
cada 1\,s, o 2PC flui do líder para as réplicas, a eleição contata os
identificadores maiores e a terminação consulta qualquer par. Não há ponto único de roteamento;
a queda de um nó não desconecta os demais. O sistema permite que cada nó
execute em uma máquina da rede local; nos experimentos deste trabalho, a
malha foi executada na própria máquina, com processos em portas distintas
(\texttt{50051..50053}), mantendo comunicação gRPC/TCP real.

A organização do código reflete a separação entre lógica e transporte:

{\footnotesize
\begin{verbatim}
distributed-db/
 proto/database.proto  contrato gRPC
 distdb/
  lamport.py           relogio de Lamport
  store.py             armazem+WAL+locks
  coordinator.py       2PC + terminacao
  election.py          alg. do valentao
  failure_detector.py  heartbeat
  node.py              servidor gRPC (no)
  client.py            cliente (failover)
 scripts/              cluster, benchmark,
                       estresse
 tests/                unidade e integracao
\end{verbatim}}

Os módulos de \texttt{distdb/} (exceto \texttt{node.py} e
\texttt{client.py}) não dependem de rede, o que permite exercitá-los em
testes de unidade simulando falhas arbitrárias.

\section{Comunicação RPC}

Toda a comunicação, tanto cliente--nó quanto nó--nó, usa gRPC. O contrato em
\texttt{database.proto} agrupa as mensagens em quatro famílias: API do
cliente (\texttt{Put}, \texttt{Update}, \texttt{Delete}, \texttt{Get},
\texttt{Exists}, \texttt{ListKeys}, \texttt{Size}, \texttt{GetAll});
transação distribuída (\texttt{Prepare}, \texttt{Commit}, \texttt{Abort},
\texttt{QueryDecision}); coordenação (\texttt{Heartbeat},
\texttt{Election}, \texttt{Announce}, \texttt{WhoIsLeader}); e recuperação
(\texttt{SyncState}). Toda mensagem carrega o relógio lógico do remetente,
e todas as chamadas usam \emph{timeouts} curtos, de modo que falhas de
comunicação sejam detectadas rapidamente em vez de bloquearem o sistema.

\section{Relógios Lógicos de Lamport}

Cada processo mantém um contador inteiro seguindo as duas regras de
Lamport~[4]: incrementa antes de todo evento local ou envio de mensagem e,
ao receber uma mensagem com carimbo $t$, atualiza para
$\max(\mathit{local},\,t)+1$. Os carimbos estabelecem a relação
\emph{happened-before} entre eventos, datam as transações no registro
durável e têm um papel central no controle de réplicas: o carimbo da última
transação confirmada por um nó (\texttt{last\_commit\_ts}) é usado para
comparar quão atualizado está o estado de cada réplica durante uma eleição
(Seção~\ref{sec:eleicao}).

\section{Exclusão Mútua Distribuída}
\label{sec:mutex}

O controle de concorrência opera em dois níveis complementares. No
\textbf{líder}, escritas concorrentes sobre a mesma chave são serializadas
por um conjunto de \emph{locks} locais (\emph{striped locks}): requisições
simultâneas de clientes diferentes são enfileiradas, enquanto chaves
diferentes seguem em paralelo. É isso que sustenta alta concorrência sem
desperdiçar trabalho com abortos (Seção~\ref{sec:aval}). Em \textbf{cada
participante}, durante a fase de votação do 2PC, o nó adquire um
\emph{lock} sobre a chave em seu armazém local; enquanto uma transação está
preparada, qualquer outra que dispute a mesma chave recebe voto NÃO e é
abortada. Esse segundo nível é a garantia \emph{distribuída}: protege a
consistência mesmo nas situações que escapam à serialização do líder, como
trocas de liderança no meio de uma transação.

\section{Eleição de Líder}
\label{sec:eleicao}

A eleição usa o algoritmo do valentão~[2]: quando um nó percebe que o líder
não responde, envia \texttt{ELECTION} a todos os nós de identificador
maior; se ninguém responde, anuncia-se coordenador para todos
(\texttt{Announce}); se algum nó maior responde, este assume a condução da
eleição. O nó vivo de maior identificador sempre vence.

Uma salvaguarda complementa o algoritmo: um nó que ficou um período fora do
ar (por exemplo, o antigo líder, justamente o de maior identificador)
poderia vencer a eleição carregando estado defasado. Por isso o vencedor,
\emph{antes} de se anunciar, consulta os pares vivos (\texttt{SyncState}) e
compara o \texttt{last\_commit\_ts} de cada um com o seu; se algum par tem
estado mais recente, o vencedor o adota. Assim a liderança nunca regride o
conteúdo do banco.

\section{Transações Distribuídas (2PC)}
\label{sec:2pc}

Cada escrita é uma transação distribuída coordenada pelo líder em duas
fases~[3]. Antes de iniciar, o líder verifica o \textbf{quórum}: se o
\emph{cohort} (ele próprio mais as réplicas vivas) é menor que a maioria do
agrupamento ($\lfloor N/2\rfloor+1$), a transação é recusada com
\texttt{no quorum} (Seção~\ref{sec:replicas}).

Na fase de \textbf{votação}, o coordenador envia \texttt{PREPARE} com a
operação a todos os participantes. Cada participante: adquire o
\emph{lock} da chave (Seção~\ref{sec:mutex}); valida a operação
(\texttt{update} e \texttt{delete} exigem chave existente); torna a
intenção durável no WAL; e vota SIM ou NÃO. Na fase de \textbf{decisão},
se todos votaram SIM o coordenador envia \texttt{COMMIT} e a escrita se
torna visível atomicamente em todos; um único voto NÃO, ou a falha de um
participante, gera \texttt{ABORT} em todos: ninguém aplica nada e os
\emph{locks} são liberados. O ponto fraco clássico do 2PC, o bloqueio de
participantes quando o coordenador morre no meio do protocolo, é tratado
pelo protocolo de terminação da Seção~\ref{sec:term}.

\section{Tratamento de Falhas}
\label{sec:falhas}

\subsection{Detecção por heartbeat}

Cada nó envia \texttt{Heartbeat} aos demais a cada 1\,s; após 3\,s de
silêncio (ou imediatamente após qualquer RPC falhar) o par é considerado
suspeito. O detector alimenta duas decisões: o líder monta o \emph{cohort}
do 2PC apenas com réplicas vivas e as réplicas percebem a morte do líder.

\subsection{Queda de réplica durante a escrita}

Se uma réplica morre no meio de uma transação, o RPC falha e a transação é
abortada (atomicidade preservada). O líder marca o nó como morto e
\emph{reexecuta} a escrita sobre o \emph{cohort} restante, de forma
transparente ao cliente, desde que haja quórum.

\subsection{Queda do líder}

As réplicas detectam a ausência de \emph{heartbeat} e disparam a eleição
(Seção~\ref{sec:eleicao}). O cliente redescobre o líder automaticamente:
ao receber \texttt{NOT\_LEADER} segue o redirecionamento e, diante de
falhas de conexão, tenta os demais nós (\emph{failover}).

\subsection{Durabilidade e recuperação}

Cada nó mantém um registro de escrita antecipada (WAL): os registros
\texttt{PREPARE}, \texttt{COMMIT} e \texttt{ABORT} são forçados ao disco
\emph{antes} de qualquer alteração em memória~[3]. Ao reiniciar, o nó
reconstrói o estado a partir do \emph{snapshot} mais recente e reaplica do
WAL apenas as transações com \texttt{COMMIT} registrado; transações
preparadas sem decisão são descartadas com segurança. Os \emph{snapshots}
compactam o log e são gerados nas transferências de estado (ressincronização
de réplica ou adoção do estado mais recente na eleição).

\subsection{Reingresso de réplica}

Ao voltar, a réplica recupera o estado local pelo WAL e pede ao líder o
estado completo (\texttt{SyncState}), ficando em dia inclusive com as
escritas que perdeu enquanto esteve fora.

\subsection{Transações em dúvida (terminação)}
\label{sec:term}

Se o coordenador morre entre o \texttt{PREPARE} e a decisão, o participante
que votou SIM fica com a chave travada sem poder decidir sozinho.
Implementamos a terminação cooperativa: transações preparadas há mais de
8\,s sem decisão fazem o nó perguntar aos pares (\texttt{QueryDecision}) o
que sabem, já que cada nó guarda de forma durável o desfecho das
transações que decidiu. Se algum par registrou \texttt{COMMIT}, o nó confirma também; se
ninguém viu decisão, aborta por presunção (\emph{presumed abort}), o que é
seguro: o coordenador só envia \texttt{COMMIT} após todos os votos SIM,
logo, se nenhum par alcançável confirmou, nenhum cliente recebeu resposta
de sucesso. Em ambos os casos o \emph{lock} é liberado e o sistema se
destrava sem intervenção manual.

\section{Controle de Réplicas}
\label{sec:replicas}

O conjunto de réplicas é gerenciado por quatro mecanismos que atuam em
conjunto: (i)~o \emph{cohort} de cada transação é montado dinamicamente com
as réplicas vivas segundo o detector de falhas; (ii)~escritas exigem
\textbf{quórum majoritário} ($\lfloor N/2\rfloor+1$): numa partição de
rede, apenas o lado com a maioria dos nós aceita escritas, e um líder
isolado responde \texttt{no quorum} (as leituras do último estado
confirmado continuam disponíveis), eliminando o \emph{split-brain} com
líderes confirmando escritas divergentes; (iii)~réplicas que reingressam
são ressincronizadas por transferência de estado; e (iv)~o vencedor de uma
eleição adota o estado mais recente entre os pares vivos
(\texttt{last\_commit\_ts}). O resultado é que o grupo de réplicas evolui
sem perder escritas confirmadas nem ressuscitar estado antigo, numa
escolha explícita de consistência em vez de disponibilidade.

\section{Avaliação Experimental}
\label{sec:aval}

A avaliação combinou três métodos. Os experimentos foram executados em um
notebook com processador Intel Core i7-1260P (12\textsuperscript{a} geração,
12 núcleos/16 \emph{threads}), 16\,GB de RAM e SSD NVMe, sob Windows~11; os
três nós e os clientes executaram na mesma máquina, comunicando-se por
gRPC/TCP pela interface de \emph{loopback}. A execução com cada nó em uma
máquina física distinta da rede local está implementada, com topologia
configurável via \texttt{--peers}, mas não foi testada nesta entrega: todos
os experimentos relatados a seguir usaram uma única máquina.

\textbf{Testes automatizados.} Suíte de unidade e integração (diretório
\texttt{tests/}) que exercita armazém, 2PC, eleição, detector, terminação e
quórum sem rede. Os cenários incluem: queda de participante na fase de
votação; coordenador que morre após confirmar em apenas um participante
(resolvido pela terminação); recuperação por WAL que reaplica somente
transações confirmadas; escrita recusada sem quórum; e o fluxo completo
escrita $\rightarrow$ queda de réplica $\rightarrow$ queda do líder
$\rightarrow$ eleição $\rightarrow$ ressincronização.

\textbf{Injeção manual de falhas.} Com o agrupamento no ar, derrubamos
réplicas e o líder em momentos distintos e acompanhamos os registros. Na
queda do líder (nó~3), o nó~2 detectou a ausência de \emph{heartbeat},
venceu a eleição e assumiu em poucos segundos (a detecção é limitada
pelo \emph{timeout} de 3\,s, mas costuma ser mais rápida, pois o RPC ao
líder falha de imediato); a réplica restante ressincronizou o estado
completo do novo líder em seguida. Também observamos a reexecução
transparente de escritas após queda de réplica, a ressincronização no
reingresso e a recusa de escritas (\texttt{no quorum}) quando restou
apenas um nó.

\textbf{Desempenho.} O \emph{benchmark} (200 operações de cada tipo)
compara o agrupamento com um armazém local único, sem rede. O armazém
local atingiu $\approx$472~mil escritas/s e $\approx$3{,}6~milhões de
leituras/s; o agrupamento sustentou 173{,}3 escritas/s (5{,}77~ms por
operação) e 2{,}7~mil leituras/s (0{,}37~ms). A diferença de mais de três
ordens de grandeza nas escritas (cerca de 2{.}700 vezes) é o preço da
durabilidade e da replicação síncrona: cada escrita percorre duas rodadas
de RPC (votação e decisão) e força o WAL ao disco (\emph{fsync}) em cada um
dos três participantes, enquanto o \emph{baseline} apenas atualiza um
dicionário em memória. As leituras distribuídas pagam somente uma RPC
local, por isso ficam cerca de quinze vezes mais rápidas que as escritas.
A Tabela~\ref{tab:perf} resume.

\textbf{Alta concorrência.} O teste de estresse disparou 8 clientes
simultâneos com 100 operações cada (800 no total; 30\% de leituras e
30\% das escritas disputando 4 chaves ``quentes''). O agrupamento
sustentou 775{,}5 ops/s agregadas e concluiu \textbf{560 escritas
confirmadas, sem nenhum aborto e nenhuma falha}: a serialização por chave
no líder enfileira as escritas concorrentes na mesma chave em vez de
deixá-las se abortarem no 2PC, ao custo de latência sob disputa
(p50 de 12{,}2\,ms, p95 de 21{,}8\,ms e p99 de 29{,}5\,ms nas escritas,
contra p50 de 2{,}2\,ms nas leituras). Ao final, o verificador leu o estado
completo de cada nó e confirmou os três \textbf{idênticos}, com 594
chaves. Repetimos o teste derrubando o líder sob carga: a vazão cai durante
a janela de reeleição (poucos segundos), os clientes reexecutam as escritas
pendentes no novo líder e nenhuma operação se perde.

\begin{table}[ht]
\centering
\caption{Resultados de desempenho (medições do grupo).}
\label{tab:perf}
\small
\begin{tabular}{lrr}
\hline
\textbf{Cenário} & \textbf{Escritas} & \textbf{Leituras}\\
 & \textbf{(ops/s)} & \textbf{(ops/s)}\\
\hline
Máquina única (sem rede)          & 471\,698 & 3\,636\,359\\
Distribuído, 3 nós (2PC)          & 173{,}3  & 2\,688{,}7\\
Estresse: 8 clientes (carga mista) & 542{,}8 & 232{,}7\\
\hline
\end{tabular}
\end{table}

\section{Desafios e Problemas Encontrados}

Os principais desafios foram: (i)~reproduzir falhas de forma controlada
para testar, resolvido separando a lógica distribuída do transporte
gRPC, o que permite simular quedas em testes de unidade, complementado por
injeção manual de falhas no agrupamento real; (ii)~o caso do coordenador
que morre no meio do 2PC, que exigiu estudar e implementar o protocolo de
terminação e a semântica de \emph{presumed abort}; (iii)~impedir que um
líder antigo reingressasse impondo estado defasado ao grupo, resolvido com
a comparação de \texttt{last\_commit\_ts} na eleição; (iv)~equilibrar
exclusão mútua e vazão sob alta concorrência, que levou ao desenho em dois
níveis da Seção~\ref{sec:mutex}; e (v)~interpretar as medições com os três
nós na mesma máquina: os processos disputam os mesmos núcleos e o mesmo
disco, e a persistência (\emph{fsync} do WAL em cada participante) domina a
latência das escritas, o que exige cuidado ao extrapolar os números para um
ambiente com máquinas separadas.

\section{Conclusão e Melhorias Futuras}

O sistema desenvolvido mantém réplicas consistentes sob falhas de
réplicas, falha do coordenador e concorrência intensa, combinando 2PC com
terminação cooperativa, eleição de líder com adoção de estado mais
recente, quórum majoritário, WAL e transferência de estado. Os
experimentos evidenciam o custo inerente da replicação síncrona nas
escritas e a recuperação automática do serviço em poucos segundos após
falhas, sem perda de operações confirmadas.

Como melhorias futuras destacam-se: repetir a avaliação com cada nó em uma
máquina física distinta da rede local (o suporte já está implementado,
com topologia configurável via \texttt{--peers}, mas os experimentos
desta entrega usaram uma única máquina); adotar um protocolo de consenso completo
(Raft ou Paxos), com garantias formais também sob partições arbitrárias;
WAL binário com soma de verificação; níveis de consistência configuráveis
nas leituras; e particionamento do espaço de chaves (\emph{sharding}) para
escalar as escritas horizontalmente.

\section*{Declarations}

\subsection*{Authors' Contributions}

{\small
Luís Augusto Starling Toledo foi responsável pela arquitetura do sistema,
pela implementação do nó (servidor gRPC) e pela integração dos
módulos.

Edson Pimenta de Almeida implementou o armazém chave--valor com WAL,
\emph{snapshots} e recuperação após falhas, além da organização do
ambiente e do repositório.

Felipe Carvalho de Paula Silva foi responsável pelo contrato Protocol
Buffers (mensagens de transação, eleição, terminação e recuperação), pela
geração dos \emph{stubs} e pelo cliente com redirecionamento e
\emph{failover}.

João Lucas de Melo Quintão desenvolveu a suíte de testes de unidade e
integração e os roteiros de injeção de falhas.

Juan Pablo Ramos de Oliveira foi responsável pela coordenação das
escritas via 2PC (Two-Phase Commit).

Luiz Gabriel Milione Assis implementou o detector de falhas por
\emph{heartbeat} e a eleição de líder, e preparou a execução do sistema em
múltiplas máquinas (topologia configurável via \texttt{--peers}).

Túlio Gomes Braga conduziu a avaliação experimental (\emph{benchmark} e
teste de estresse), com análise dos registros e dos resultados.
\par}

\section*{References}

{\small
\noindent [1] Google. gRPC Documentation. Disponível em:
\url{https://grpc.io/}

\noindent [2] Garcia-Molina, H. (1982). Elections in a distributed
computing system. \emph{IEEE Transactions on Computers}, C-31(1):48--59.

\noindent [3] Gray, J. e Reuter, A. (1993). \emph{Transaction Processing:
Concepts and Techniques}. Morgan Kaufmann.

\noindent [4] Lamport, L. (1978). Time, clocks, and the ordering of events
in a distributed system. \emph{Communications of the ACM},
21(7):558--565.

\noindent [5] Tanenbaum, A. e van Steen, M. \emph{Distributed Systems:
Principles and Paradigms}. Prentice Hall.
\par}

\end{document}
